\subsection{What the evidence supports}\label{subsec:what-the-evidence-supports}

The degeneracy hypothesis --- $\varphi_c$ and $\psi$ carry no strain-only recoverable signal for this population in this architecture/loss family --- is the best-supported reading of nine training campaigns:
\begin{enumerate}
\item validation circular loss flat at the random value $\approx 1.0$ in every run, every model, every $\lambda$, with healthy gradients and unit-scale magnitudes over the second half of training in the clean models;
\item 0 of 8 $\varphi_c/\psi$ bootstrap tests significant (11 of 12 overall), with the single scalar-head detection explained as a population bias;
\item no SNR-dependent improvement anywhere --- the direction any real strain-derived effect must point;
\item the combination-space model, purpose-built to exploit the analytically confirmed conditioning structure of \S\ref{sec:phicpsi:prereq}, collapsing \emph{into} the constant predictor exactly as it should if the residual per-sample information is nil;
\item positive controls learned to high fidelity from the same features throughout;
\item no model recovers an edge-on signal exceeding its own uninformative $\iota$-control's sampling noise (\S\ref{sec:phicpsi:inclination}) --- the one population axis along which the analytic study predicts the degeneracy should weaken shows no cross-model, edge-on-favoring trend either.
\end{enumerate}

Two points sharpen what is and is not being claimed.
First, the claim is \emph{conditional}: it applies to point-estimating regression on this SNR-7--15, two-detector, dominant-mode population, under circular loss --- not to Bayesian posterior recovery (a sampler still recovers a meaningful \emph{joint} posterior over $(\varphi_c, \psi)$), not to louder signals, not to alternative data representations (e.g.\ frequency-domain or time--frequency inputs), and not to configurations with inclination supplied as side information (\S\ref{sec:phicpsi:future}).
Second, the analytic study of \S\ref{sec:phicpsi:prereq} makes the null \emph{informative} rather than merely disappointing: a combination-conditioning ratio of only ${\sim}1.16\times$ population-averaged --- concentrated face-on, where the curriculum then down-weights it to zero --- quantifies how little marginal structure there was for a per-event point estimator to find.
The networks did not miss an obvious target; they confirmed, at scale, how small the target is.

\subsection{Relation to full-posterior methods}
\label{sec:phicpsi:npe}

A natural objection to testing $\varphi_c/\psi$ recoverability under point estimation at all is that the field's standard answer to exactly this kind of degenerate parameter has existed for years: neural posterior estimation.
Conditional variational autoencoders~\citep{gabbard2022vitamin}, direct neural approximations to the posterior density~\citep{chua2020learning,green2020complete}, and, more recently, conditional normalizing flows~\citep{papamakarios2021normalizing} trained by simulation-based inference~\citep{cranmer2020frontier,dax2021dingo,dax2023dingois} do not attempt to output a single value for a parameter the likelihood constrains only partially; they instead learn an amortized approximation to the full Bayesian posterior, and demonstrably reproduce multimodal, disconnected posterior structure --- including, in the importance-sampling-corrected variant, an asymptotically exact correction against the true likelihood~\citep{dax2023dingois}.
This is not a single team's result: independent groups converged on the same amortized-posterior strategy for exactly this class of degenerate, multimodal gravitational-wave parameters~\citep{chua2020learning,green2020complete,gabbard2022vitamin,dax2021dingo}, which is why we treat it as the field's settled answer rather than one candidate approach among several.
For a genuinely degenerate target such as $\varphi_c/\psi$ at low inclination, a posterior-estimating network faces no obstacle in principle: it can represent the ridge or the multimodal structure directly, exactly as a stochastic sampler does~\citep{veitch2015lalinference,ashton2019bilby}.
In this sense the underlying scientific question --- what does the strain constrain about $\varphi_c$ and $\psi$ --- is not open; it is answered by the analytic degeneracy structure of \S\ref{sec:phicpsi:prereq} and represented correctly by any calibrated full-posterior method.

What this paper tests is narrower and, we argue, still useful: whether a \emph{point-estimating} regression head --- the architecture used for every other parameter in this study, and a component still found in fast detection-and-characterization pipelines that were not designed as full posterior samplers~\citep{george2018deep} --- fails safely or silently on a target of this kind.
The answer is silently, for a long stretch of the diagnostic arc (\S\ref{sec:phicpsi:arc}): two genuine optimization pathologies produced a plausible-looking $R^2 = 0.754$ and non-trivial gradient flow before either was traced to its mechanism, and only a pre-registered verification battery closed the question convincingly.
The contribution here is therefore not a competing answer to what full-posterior methods already answer correctly, but a documented account of how a point-estimation pipeline behaves --- and how easily it can be misread --- when it is pointed at a target the field already knows to be degenerate, together with the diagnostic discipline (\S\ref{subsec:methodological-lessons}) that a practitioner building or auditing such a pipeline can reuse.

\subsection{Sufficiency of the architecture pool}
\label{sec:phicpsi:pool}

A natural objection to any architecture-crossed null is: \emph{why only seven models?}
The answer is that architecture count is the wrong axis for defending this particular claim, and the design reflects that deliberately.

A comparative claim (``architecture $X$ is best for task $Y$'') grows stronger with breadth, and every untried model is a hole in it.
A null of the form ``the input--label relationship carries no recoverable signal'' does not: it grows stronger with \emph{diversity of inductive bias} and with \emph{mechanism-level evidence}, both of which saturate quickly.
On the first, the five trunk families (\S\ref{sec:phicpsi:trunks}, Table~\ref{tab:phicpsi:trunks}) span qualitatively different hypotheses about how phase information could be encoded --- long-range dilated convolution, local hierarchical convolution, content-based attention, explicit multi-scale filtering, and deep residual composition --- rather than five samples from one family.
On the second, the decisive observation is that \textbf{capacity was demonstrably not the binding constraint}: the higher-capacity trunks drove \emph{training} circular loss down to ${\approx}0.49$--$0.60$ while validation loss stayed pinned at 1.0 (Fig.~\ref{fig:phicpsi:comboloss}).
That is precisely what added capacity does against an uninformative target --- it memorizes.
An eighth architecture buys another memorization curve, not information absent from the strain; once the failure is localized in the data by flat validation loss under certified-healthy gradients, in agreement with the analytic degeneracy of \S\ref{sec:phicpsi:prereq}, architecture search has no remaining expected payoff.
In one sentence: \emph{the memorization gap shows the null is a statement about the data, not about the trunk architecture.}
For the angular targets specifically, this claim is scoped to the trunk and does not extend to the periodic-head design: as \S\ref{sec:phicpsi:inclination} details, $\iota$, $\varphi_c$, and $\psi$ share the same two-layer MLP head, so a head-capacity limitation specific to that shared design remains a residual, untested confound for the angular parameters that the scalar-control evidence above does not cover.

This memorization-gap argument was illustrated in aggregate (Fig.~\ref{fig:phicpsi:comboloss}, all seven architectures); an adversarial reviewer correctly asked whether it holds for the two models the null is actually \emph{certified} on, individually, rather than for the pool as a whole.
It does, but asymmetrically, and the asymmetry is worth reporting rather than folding into the aggregate.
cnn\_attention's own training/validation circular loss (Fig.~\ref{fig:phicpsi:certifiedmemorization}) reproduces the pool-level pattern exactly: training loss falls to $0.60$ ($\varphi_c$) and $0.52$ ($\psi$) by epoch 79 while validation stays at $1.01$/$1.00$ --- this model demonstrably has, and uses, spare capacity to fit sample-specific structure, and still generalizes none of it.
poc\_b does not: its combo-A/combo-B training loss tracks validation almost exactly throughout training ($0.988$/$0.998$ vs $0.999$/$0.991$ at epoch 79), the same flat-everywhere pattern shown by the other two TCN-trunk configurations, poc\_a and tcn (train $0.98$--$0.97$ vs val $1.02$--$1.01$ at epoch 79), regardless of head parameterization.
This is not evidence that the TCN trunk lacks capacity: the identical trunk reaches chirp-mass $R^2 = 0.97$ (train) / $0.96$ (val) on poc\_a --- ample nonlinear capacity when a learnable target exists.
The two certified models fail for different reasons, and the two readings do not carry equal evidential weight --- they are kept separate below rather than read together.
This section's architecture-sufficiency argument rests on cnn\_attention's memorization gap alone: it finds and fits sample-specific structure it cannot generalize (the overfitting-without-learning pattern already qualified in \S\ref{sec:phicpsi:run7}), the cleanest single demonstration that spare capacity finds and cannot generalize structure from an uninformative target.
poc\_b's optimizer finds no structure even in-sample, on the same TCN trunk family that reaches chirp-mass $R^2 = 0.97$ elsewhere --- but this pattern is not a second, independent instance of the same check.
The curriculum weight $w(\iota) = \sin^2\iota$ suppresses the poorly-constrained combination's gradient for near-face-on samples ($w < 0.19$ for the $28.7\%$ of the population with $|\cos\iota| > 0.9$), so a flat in-sample loss under this curriculum is equally consistent with an unlearnable target and with curriculum-induced gradient starvation over most of the dataset.
poc\_b (\S\ref{sec:phicpsi:confounds}) therefore remains a distinct, curriculum-confounded consistency check that corroborates from the optimizer side without carrying the same evidential weight as cnn\_attention's memorization gap --- not a second capacity-sufficiency certification.

\begin{figure}[tbp]
    \centering
    \includegraphics[width=\textwidth]{certified_models_train_val_loss.png}
    \caption[Train vs.\ validation circular loss, certified models only]{Train vs.\ validation circular loss for the two certified models specifically, generated directly from each run's \texttt{history.csv} (\texttt{plot\_certified\_memorization.py}, no model loading or GPU use). cnn\_attention (right column) shows the memorization gap (train $\downarrow$ to $0.60$/$0.52$, val flat $\approx 1.0$); poc\_b (left column) shows no gap at all --- train tracks validation throughout, the same pattern shown by the other two TCN-trunk configurations (poc\_a, tcn).}
    \label{fig:phicpsi:certifiedmemorization}
\end{figure}

The progressive narrowing of the pool --- seven configurations, then four, then two --- likewise reflects criteria, not attrition.
The seven were never seven independent tests but five trunks plus the two head parameterizations on the primary trunk; the cut to four (poc\_a, poc\_b, tcn, cnn\_attention) selected the $\lambda$-matched subset after the magnitude penalty was introduced (\S\ref{sec:phicpsi:protocol}), the remaining three trunks retaining evidential value as unmatched corroboration; and the cut to the two \emph{certified-clean} models (poc\_b, cnn\_attention) was imposed by the pre-declared std\_ratio interpretability gate (\S\ref{sec:phicpsi:confounds}), not by preference among results.
Each restriction is documented and criterion-based, and the excluded configurations agree with the included ones at every point where they can be compared.

For the same reason, the deliberate decision \emph{not} to expand the architecture pool in future work (\S\ref{sec:phicpsi:future}) is scoping rather than omission: the evidence localizes the failure in the information available to the network, so the productive next lever is conditioning information ($\iota$-conditioning), not a broader model zoo.
A single modern-family spot-check (e.g.\ a state-space sequence model at $\lambda = 0.01$) would be a cheap robustness supplement, but per unit of compute, seed replication of an existing configuration addresses a larger residual risk (\S\ref{sec:phicpsi:threats}) than any additional architecture.

\subsection{Threats to validity}
\label{sec:phicpsi:threats}

We record the limitations explicitly, in roughly descending order of concern.
\begin{itemize}
\item \textbf{Two configurations remain uninterpretable.} tcn/$\varphi_c$ and poc\_a/$\psi$ never achieved certified metric health at any $\lambda$; the clean degeneracy test rests on poc\_b and cnn\_attention (plus the ablation and pre/post-fix consistency of every other configuration).
The null is multiply corroborated, but its \emph{certified} base is narrower than its table count.
\item \textbf{The A.3 closure rests on a within-run-validated channel.}
The trace's pre-stated displacement-geometry classifier failed its positive control at calibration and was retired (\S\ref{sec:phicpsi:asymmetry}); the closure rests on the paired probe-loss statistic, whose own control passed emphatically in the same run (early mchirp $|t| = 3.4$--$8.5$) but which was added on review rather than pre-registered.
The evidence is positive-controlled; the channel choice is post hoc --- recorded here rather than hidden.
\item \textbf{Waveform-model provenance is a reproducibility gap, not merely a caveat.}
All injections are understood to use the dominant quadrupole ($\ell=2$, $|m|=2$) mode only (IMRPhenomD, no higher-order modes) --- the condition under which the face-on degeneracy argument of \S\ref{sec:phicpsi:intro} is exact --- but this is not preserved as metadata in the HDF5 dataset, nor in any generation script retained in the accompanying repository; it is recorded here from the authors' knowledge of how the dataset was produced.
The primary concern this raises is reproducibility, not the direction of the null: a reader cannot regenerate or independently verify the dataset from what is provided, and on review the honest position is that this is not acceptable practice, not a footnote --- any follow-up work should regenerate the dataset from a version-controlled script that records its own configuration before relying further on this result.
Separately, and only as a secondary point, the direction of the risk is asymmetric: if this recollection were wrong and higher-order modes were in fact present, the degeneracy would be only approximately exact even face-on, which would \emph{sharpen} rather than weaken the null (a network failing to exploit even a HOM-broken degeneracy is stronger evidence of no recoverable signal) --- but this does not excuse the provenance gap, which is a threat to the paper's reproducibility independent of which way it would cut if resolved.
\item \textbf{Residual anomalies.} tcn/$\psi$'s $+0.0072$ validation-loss drift at $\lambda = 0$ remains unexplained (small, on a loss of ${\sim}1.0$); the inclination head's failure mechanism (Huber path, no normalization) is documented as distinct but is itself unresolved; and a sky-position degradation specific to the \texttt{SumDiffTrainer} runs (\SIrange{8.2}{10.0}{\degree} vs \SIrange{3.3}{4.5}{\degree} for the plain baselines) was flagged in the record but never investigated --- a known open issue, out of scope for this paper.
\item \textbf{Design constraints.}
Single seed (42) per configuration --- replication is across architectures and $\lambda$ values, not across seeds; the five-trunk comparison is not $\lambda$-matched (three trunks trained at $\lambda = 0$), though the four models carrying the verified null are matched; the curriculum used the analytic $\sin^2\iota$ weight, not the Jacobian-fitted variant it was validated against; fixed 80-epoch budget with no early stopping; a single simulated dataset with a single noise realization per event; and the inclination prior is uniform in $\iota$ itself rather than in $\cos\iota$, the astrophysically correct isotropic choice, which over-represents edge-on systems relative to reality --- since edge-on is where the degeneracy is weakest, this deviation does not appear to bias the claim in a self-serving direction, but it should have been isotropic and was not.
\item \textbf{Metric-versioning note.}
Two evaluation passes over the Run~7 checkpoints differ by up to \SI{0.026}{\radian} in periodic ang\_MAE (both straddling the same null; no verdict affected).
Tables~\ref{tab:phicpsi:prepost} and~\ref{tab:phicpsi:bootstrap} use the self-consistent bootstrap/stratification set.
\end{itemize}
None of these, alone or jointly, provides a mechanism by which a real, useful $\varphi_c/\psi$ signal could have hidden from all of \S\ref{sec:phicpsi:results} --- but they bound the strength of the claim honestly.

\subsection{Methodological lessons}\label{subsec:methodological-lessons}

Three transferable lessons emerge from the diagnostic arc, and they are the part of this paper most likely to outlive its specific numbers, echoing a caution already raised elsewhere about over-reading aggregate performance claims for deep learning applied to gravitational-wave data~\citep{gebhard2019magicbullet}.

\paragraph{Aggregate metrics are not evidence until their mechanism is verified.}
Every major error in this investigation was a plausible reading of a true number: $R^2 = 0.75$ (mode collapse in disguise), an endpoint std\_ratio of 0.62 (concealing a crash to 0.07 and recovery), rising validation loss (a regularizer interaction), a ``significant'' $p = 0.0007$ (population bias).
The correctives were never better metrics but mechanism inspection: forward-pass dumps, gradient-chain traces, full trajectories, code-level loss-path audits.
For null results especially, the epistemic burden inverts: a \emph{flat} loss is only informative once the gradient path that would have moved it is certified live.

\paragraph{Controls must be validated at the code level.}
The inclination head's coincident failure was, for a time, load-bearing evidence for a shared-mechanism (hence fixable, hence non-physical) explanation.
It shared the encoding, the architecture, and the failure signature --- but not, it turned out, the loss path.
A control that resembles the treatment in every visible metric can still differ in the one line of code that matters.

\paragraph{Pre-registration has a place inside engineering loops, not just at study level.}
The $\lambda$ retune's gates were written down precisely because the investigators had watched themselves misread results three times.
The payoff was concrete: a near-miss that endpoint-eyeballing would have waved through was held to the written criterion, a mid-stream temptation to recalibrate the gate was resisted and documented instead, and the final filing --- ``neither null nor counter-evidence'' --- is a category that post-hoc reasoning almost never produces but that honest accounting sometimes requires.

\subsection{Future work}
\label{sec:phicpsi:future}

Seven successors are scoped, in rough priority order --- though (iv) requires no retraining and could be run independently of where the others land in the queue.
(Two items originally listed here as immediate checks have since run and are no longer future work: the inclination-stratified breakdown, \S\ref{sec:phicpsi:inclination}, and the scalar-control chirp-mass-by-inclination-band check, Table~\ref{tab:phicpsi:controlstrat}.)
(i)~\textbf{Inclination conditioning}: supply true $(\sin\iota, \cos\iota)$ as an auxiliary input --- the analytic structure of \S\ref{sec:phicpsi:prereq} says the well-constrained combination is knowable \emph{given} $\iota$, so this tests whether the degeneracy is breakable with side information; a full implementation plan exists (train-time truth, with inference-time $\iota$ estimation explicitly out of scope).
(ii)~An \textbf{architecture-level attack on the std\_ratio instability} for the two uninterpretable pairs (the pre-registered $\lambda$-sweep's named next lever), or an explicit decision to close that thread; alongside it, a finer freshly pre-registered $\lambda$ mini-sweep over $[0.02, 0.08]$, motivated by the peaked $\lambda$-response observed in \S\ref{sec:phicpsi:prereg}.
(iii)~A \textbf{posterior-estimation reformulation}: the natural follow-on from ``point estimation is hopeless'' is not resignation but a conditional-density head (e.g.\ a von Mises mixture over the well-constrained combination), which the present paper's null both motivates and baselines.
(iv)~A \textbf{ridge-structure check}: a 2D joint scatter of the certified models' predicted $(\hat\varphi_c, \hat\psi)$ pairs, to test directly whether the network learned the degenerate combination manifold itself (which would appear as a diagonal ridge) rather than the collapsed or spread patterns already documented in Fig.~\ref{fig:phicpsi:hist} --- raised during internal review; inference-only against saved checkpoints, no retraining required, but not run here because no script or output artifact yet backs it (the same artifact-traceability standard \S\ref{sec:phicpsi:threats} sets for the current text applies equally to future additions).
(v)~A \textbf{synthetic-data ablation for poc\_b}: retrain the combination-head design on a synthetic dataset in which the well-constrained combination is recoverable by construction, disentangling the curriculum-design confound (\S\ref{sec:phicpsi:confounds}, \S\ref{sec:phicpsi:pool}) from the physical degeneracy --- the direct test of whether poc\_b's collapse is what its design does to an unlearnable target rather than an artifact of the curriculum itself.
(vi)~A \textbf{sensitivity-floor check}: a high-SNR (25--30) validation set, or a Fisher-matrix bound~\citep{vallisneri2008fisher} on the achievable $\varphi_c/\psi$ accuracy at this population's SNR range, to rule out the possibility that the null reflects a detection-scale sensitivity floor rather than the degeneracy as such.
(vii)~\textbf{Dataset regeneration from a version-controlled generator script}, recording approximant, mode content, and priors as dataset metadata --- closing the waveform-provenance gap of \S\ref{sec:phicpsi:threats} for good and making the null independently reproducible.